\subsection{An unpaid participating dividend can change ordinary earnings}
\label{participating-tax-allocation}

Assumed conversion is not the only hypothetical transaction that changes
earnings per share. Allocating undistributed profit between ordinary and
participating shares can itself have a tax consequence. Suppose participating
dividends are tax deductible and the resulting benefit belongs to ordinary
shareholders. An allocation that divides reported profit solely by the base
participation ratio misses that additional ordinary entitlement. The June 2017
IFRS Interpretations Committee decision requires the ordinary numerator to
include the ordinary holders' share of the tax benefit from the hypothetical
distribution required by IAS~33.A14. This applies whether an actual benefit
is recognised in profit or directly in equity
\citep{ifric2017participatingtax}.

The accompanying example makes the consequence visible before any dividend
is paid. Participating and ordinary holders share base distributions in the
ratio ten to one. The tax deduction is 30\% of the participating distribution,
and its entire benefit belongs to ordinary holders. In the first year, profit
before that benefit is 330 and no dividend is declared. The ordinary base
share is \(330/11=30\), while participating holders' hypothetical allocation
is 300. Distributing that 300 would create a benefit of 90 for ordinary
holders. Their basic-EPS numerator is therefore \(30+90=120\), although no
cash dividend or actual distribution-related tax benefit arose in that year
\citep[pp.~1--5]{ifric2017participatingtaxexample}.

To carry this allocation across periods, distinguish the profit being shared
from dividends that may distribute earlier earnings. Let \(P\) be current
profit before the participating-dividend tax benefit, \(r\) the participating-
to-ordinary base-rights ratio, and \(t\) the tax rate. Define the class
fractions \(a=r/(1+r)\) and \(b=1/(1+r)\). If the base dividend declared
in the current period is \(D\), participating holders receive \(aD\) and
ordinary holders receive \(bD\). In the source's payout pattern the actual
tax benefit \(H=taD\) is also distributed to ordinary holders in that period.
All quantities here are money except the ratios and tax rate; classification,
deductibility, recognition and these rights are supplied facts.

First put the actual tax benefit on a consistent basis for the EPS allocation.
Reported profit is \(P+H\) when that benefit is recognised in profit, or
\(P\) when it is recognised directly in equity. Adding the latter benefit
for this allocation gives adjusted profit \(G=P+H\) in either case. Subtract
all declared class dividends, including the ordinary tax-benefit dividend:
\[
 R=G-aD-bD-H=P-D.
\]
This signed residual is the base profit still to allocate under the hypothetical
distribution. It can be negative when the actual dividend includes earlier
years' earnings. The participating residual is \(aR\), the ordinary base
residual is \(bR\), and the ordinary holders' further tax adjustment is
\(taR\). Adding declared and residual amounts gives
\[
 \begin{aligned}
 N_O&=bD+H+bR+taR,\\
 N_P&=aD+aR,\qquad \operatorname{EPS}_{\mathrm{basic}}=N_O/n,
 \end{aligned}
\]
where \(N_O\) and \(N_P\) are the ordinary and participating earnings
allocations and \(n>0\) is the supplied weighted-average ordinary denominator.
These are the example's three steps: adjust and subtract declared dividends,
allocate the remaining profit and its tax consequence, then add the declared
allocations back to each class
\citep[Steps~1--3, pp.~3--5]{ifric2017participatingtaxexample}.

The example's following two years show why the signed residual matters.
Base profits are 550 and 110. The second year distributes the first year's
base profit of 330; the third distributes the remaining base profit of 660.
All actual tax benefits are paid to ordinary holders. The complete allocation is
\begin{center}
\begin{tabular}{lrrr}
\toprule
 & Year 1 & Year 2 & Year 3\\
\midrule
Current base profit \(P\) & 330 & 550 & 110\\
Base dividend \(D\) & 0 & 330 & 660\\
Actual benefit \(H\) & 0 & 90 & 180\\
Residual \(R=P-D\) & 330 & 220 & \(-550\)\\
Residual tax adjustment \(taR\) & 90 & 60 & \(-150\)\\
Ordinary earnings \(N_O\) & 120 & 200 & 40\\
Participating earnings \(N_P\) & 300 & 500 & 100\\
\bottomrule
\end{tabular}
\end{center}
In year 3 the ordinary dividend is \(60+180=240\). Its remaining base
allocation is \(-50\), and the tax adjustment on that residual is \(-150\),
giving \(240-50-150=40\). Setting the negative adjustment to zero would give
190 and count part of a benefit already attributed to earlier earnings again.
The \(-150\) is an EPS-allocation reversal; it is neither a negative cash
dividend nor evidence that every tax loss creates a recognised asset.

The dependence on payout timing cancels under these fixed rights. Substituting
\(H=taD\) and \(R=P-D\) into the class allocations yields the local identities
\[
 \begin{aligned}
 N_O&=bD+taD+(b+ta)(P-D)=(b+ta)P,\\
 N_P&=aD+a(P-D)=aP,\\
 N_O+N_P&=P+taP=G+taR.
 \end{aligned}
\]
Thus current ordinary earnings depend on current profit and its associated
hypothetical benefit, even when the declared dividend distributes older profit.
At zero tax the formula reduces to the ordinary base share \(bP\); with no
participating entitlement it reduces to \(P\). Total class allocations can
exceed reported profit because they include the hypothetical tax consequence.
For example, year 1 allocates \(120+300=420\) while reported profit is 330.
Those two quantities serve different purposes and must not be forced to match
by posting the hypothetical difference.

\texttt{ParticipatingTaxEPS/1} specifies this basic-EPS calculation with exact
rational amounts and a supplied denominator. It checks actual class dividends,
the benefit amount and recognition location, and an opening base-profit balance
sufficient to fund any prior-period distribution. It permits a negative
post-dividend residual while restricting current base profit to nonnegative
amounts. Missing rights or tax findings leave the result unresolved;
inconsistent payouts, future information and unsupported findings produce no
partial numerator. Its output contains no dividend, tax asset or journal entry.

The interpretation also says that a tax benefit shared by both classes must
follow the applicable rights \citep[p.~2, footnote~(a)]{ifric2017participatingtaxexample}.
That is a different allocation from the ordinary-only benefit used here.
Retained actual benefits, changing or capped rights, loss-tax recognition and
nonadditive diluted selection require their own supported treatment. In
particular, a fixed after-tax earnings input to the earlier capped-participation
method cannot silently stand in for the hypothetical benefit calculated here.
